\subsection{An isolated dyadic wave and an actual-coefficient index test}

The simple poles of P-F31 permit a different strict-index test from
the average-polynomial test in P-F34. This one uses a specified
cyclotomic trace and an elementary bound on an actual coefficient.
No sign or lower bound for that trace is asserted.

\begin{proposition}[P-F56: maximal two-part poles and paired constituents]
\label{prop:cubicpairedwave}
Use the ordinary cubic multiplicity algebra $T^{(n)}$ and its
Hilbert function $f_n$ from P-F31, with $n\ge4$, and put
\[
 M=2^{n-1},\qquad g=\gcd(n,3),\qquad d_0=n/g,
 \qquad m=\binom{n+2}{3}.
\]
Every pole of $H_n(t)$ is a root of unity whose order has
two-part at most $M$. Among roots whose order has two-part exactly
$M$, the only poles are the primitive $M$th roots, and each of
those poles is simple. More precisely, if $\eta$ is a root of
unity of order $q$ with $v_2(q)\ge n-1$, then
\[
 \eta\text{ is a pole of }H_n
       \quad\Longleftrightarrow\quad q=M;
 \qquad q=M\Longrightarrow\text{the pole is simple}.
 \tag{P-F56a}
\]

Let $L$ be the least quasipolynomial period of $f_n$, and write
$f_n(s)=q_r(s)$ for $s\equiv r\pmod L$, with
$0\le r<L$. Then
\[
                     v_2(L)=n-1.                            \tag{P-F56b}
\]
Fix $\zeta=e^{2\pi i/M}$ and retain the actual-residue
normalization of P-F34:
\[
 \mathcal C_n(\zeta)=\lim_{t\to\zeta}(1-t^M)H_n(t),
 \qquad A_n(\zeta)=\mathcal C_n(\zeta)/M.
\]
For each integer $r$ define the rational, $M$-periodic wave
\[
 \begin{split}
 W_n(r)
   &=\sum_{\substack{1\le a<M\\a\text{ odd}}}
                      A_n(\zeta^a)\zeta^{-ar}\\
   &=\frac1M\operatorname{Tr}_{\mathbb Q(\zeta)/\mathbb Q}
                  \bigl(\mathcal C_n(\zeta)\zeta^{-r}\bigr).
 \end{split}                                                 \tag{P-F56c}
\]
For every residue $r$ and every complex argument $X$ one has
the exact polynomial identity
\[
               q_r(X)-q_{r+L/2}(X)=2W_n(r),                 \tag{P-F56d}
\]
where subscripts are read modulo $L$. In particular,
\[
 \begin{split}
 f_n\bigl(p(f_n)\bigr)&\ge2W_n\bigl(p(f_n)\bigr),\\
 f_n\bigl(s(f_n)+1\bigr)&\ge2W_n\bigl(s(f_n)+1\bigr).
 \end{split}                                                 \tag{P-F56e}
\]
These are necessary conditions for the literal indices of
GCT VI, Definitions~1.2.2 and 1.2.4, printed pp.~13--14.
They require no bound on the average polynomial $B_n$ or on
Fourier terms of smaller two-part.
More generally, the inequality $f_n(x)\ge2W_n(x)$ holds
at every integer $x\ge p(f_n)$ and at every integer
$x\ge s(f_n)+1$. Thus a single integer $x\ge0$ satisfying
\[
          2W_n(x)>f_n(x)
       \quad\Longrightarrow\quad p(f_n)>x,
                                  \quad s(f_n)\ge x.
 \tag{P-F56e'}
\]

For every fixed real $K\ge0$ and every integer $0\le x\le n^K$,
\[
 f_n(x)\le\binom{m+d_0x-1}{m-1}
     \le\bigl(2n^{\max\{3,K+1\}}\bigr)^{n^3}.
 \tag{P-F56f}
\]
Consequently the following would suffice to refute both narrow
polynomial-index hypotheses on this cubic family: there exist
$\epsilon>0$ and an unbounded set $\mathcal N$ of dimensions
such that, for every fixed $K\ge0$ and every sufficiently large
$n\in\mathcal N$, there is an integer $x$ satisfying
\[
 \lceil n^K\rceil\le x\le\lceil n^K\rceil+2,
                         \qquad W_n(x)\ge2^{\epsilon n^4}.
 \tag{P-F56g}
\]
Under this additional condition, both $p(f_n)$ and $s(f_n)$
exceed every fixed polynomial in $n$ along $\mathcal N$.
Condition (P-F56g) is not proved here. The wave can have
either sign a priori, and the absolute value of one Fourier
coefficient in P-F34 does not establish (P-F56g).
The window of three consecutive integers is a sufficient
target, not a proved oscillation pattern; any bounded-width
window, or suitable polynomially bounded witnesses tending
beyond every fixed polynomial threshold, would work as well.
\end{proposition}

\begin{proof}
We first prove a fixed-point criterion for absence of a Hilbert
pole, including possible nonreduced structure at the origin.
Let $R$ be any finitely generated positively graded complex
algebra with $R_0=\mathbb C$, and let $\eta$ have order $q$.
The ideal defining its grading fixed scheme is
\[
                    I_q=(R_d:q\nmid d).                     \tag{P-F56h}
\]
Indeed a homogeneous element of degree $d$ changes by
the scalar $\eta^d$, so its difference from its transform
is zero exactly when $q\mid d$.
Choose homogeneous ideal generators $a_1,\ldots,a_l$ of
$I_q$, each of degree $e_i$ not divisible by $q$.
Such a finite choice is possible by Noetherianity, starting
with the indicated homogeneous generating set.
If the fixed scheme is supported only at the graded origin,
then $\sqrt{I_q}=R_{>0}$ and $R/I_q$ is a finite-dimensional
complex algebra. To see this directly, a power of the finitely
generated ideal $R_{>0}$ lies in $I_q$, and only finitely many
monomials of bounded length in homogeneous algebra generators
survive.

The homology modules $N_j$ of the graded Koszul complex
$K_R(a_1,\ldots,a_l)$ are finite $R$-modules annihilated by
$I_q$: multiplication by each $a_i$ is null-homotopic on
the Koszul complex. Thus all $N_j$ are finite-dimensional.
The degreewise Euler identity gives
\[
 H_R(t)\prod_{i=1}^l(1-t^{e_i})
                    =\sum_j(-1)^j\operatorname{Hilb}_{N_j}(t).
 \tag{P-F56i}
\]
The right side is a polynomial and every factor on the left
is nonzero at $t=\eta$. Hence $H_R(t)$ is regular at $\eta$.
This is the zero-dimensional case of the Koszul calculation
in P-F23; reducedness of the fixed scheme is unnecessary.

Apply this criterion to $R=T^{(n)}$ and $X=\operatorname{Spec}R$.
Let $q=2^ab$ with $b$ odd and $a\ge n-1$.
Suppose a nonzero closed point $p\in X$ is fixed by a primitive
$q$th root $\eta$. It is then fixed by
$\eta_2=\eta^b$, which has order $2^a$.
We need the closed-orbit lift for this pure two-primary
subgroup, with no assumption about the odd part of $q$.
Here are the details extending the first part of P-F26.

Put $\gamma=3/g$, an odd integer. The canonical grading
character on quotient points is
$(\det B)^{-\gamma}$ for $B\in GL_n$, as in (P-F26b).
There is a primitive $2^a$th root $\theta$ such that
$\theta^{-\gamma}=\eta_2$.
Choose $B\in GL_n$ with $\det B=\theta$, and choose
a cubic $F$ in the unique closed $SL_n$ orbit over $p$.
Both $F$ and $B\cdot F$ have closed orbits and the same
quotient point. Therefore $B\cdot F=S\cdot F$ for some
$S\in SL_n$, and $C=S^{-1}B$ fixes $F$ exactly with
$\det C=\theta$.
The unique closed orbit in each affine reductive quotient
fiber is the standard input already used in P-F26; an exact
primary locator is Luna, \emph{Slices \'etales},
Section~I.1, printed p.~84:
\url{https://www.numdam.org/item/10.24033/msmf.110.pdf}.

The semisimple part $C_s$ also fixes $F$ and has determinant
$\theta$. Indeed in the commuting Jordan decomposition,
an eigenspace of $C_s$ with eigenvalue different from one
cannot contain a fixed vector of $C_sC_u$, because the
unipotent part $C_u$ has all eigenvalues one.
Diagonalize $C_s$ and let $\Gamma\subset\mathbb C^*$ be
the finitely generated group generated by its eigenvalues.
Project $\Gamma$ to its finite two-primary torsion subgroup,
using a decomposition into its finite cyclic torsion subgroup
and a free complement. This projection fixes $\theta$.
Applying it to every eigenvalue preserves every monomial
relation expressing $C_sF=F$.
The resulting diagonal matrix $A$ fixes $F$, has finite
two-power order, and has determinant of order exactly $2^a$.

Since $p$ is nonzero, a positive-degree invariant is nonzero
at $F$. In any eigenbasis, one of its nonzero coefficient
monomials uses only monomials supported in $F$.
Invariance under the diagonal $SL_n$ torus forces the
corresponding nonzero nonnegative integral combination of supported
cubic exponent vectors to be a positive rational multiple
of $\mathbf1$. Thus $\mathbf1$ belongs to the rational
span of the support, precisely the hypothesis of P-F30.
Theorem P-F30 gives
\[
                       2^a=\operatorname{ord}(\det A)\le M.
 \tag{P-F56j}
\]
If $a>n-1$ this is a contradiction, so no nonzero fixed
point exists. If $a=n-1$, equality in P-F30 makes $F$
equivalent to the chain cubic, hence $p$ lies on its
nonzero grading orbit. The exact stabilizer of every
point of that orbit is $\mu_M$ by P-F28, including the
converse lift argument in (P-F28d).
The original root $\eta$ can fix $p$ only if $\eta\in\mu_M$.
Since its order is $q=Mb$, this requires $b=1$.

It follows that for every $q$ in (P-F56a) except $q=M$
the fixed scheme has no nonzero closed point. Its support
is consequently just the graded origin, because a closed
subset of an affine scheme of finite type over $\mathbb C$
is determined by its closed points. Equations
(P-F56h)--(P-F56i) show that there is no Hilbert pole there.
For $q=M$, P-F31 proves an actual simple pole at every
primitive root. This proves all of (P-F56a).

For completeness, the least quasipolynomial period is the
least common multiple of the orders of the Hilbert poles.
Indeed the Fourier polynomials of P-F34 are
\[
 P_\rho(X)=\frac1L\sum_{r=0}^{L-1}\rho^r q_r(X),\qquad
 q_r(X)=\sum_{\rho^L=1}\rho^{-r}P_\rho(X).                 \tag{P-F56k}
\]
The generating function of a nonzero polynomial times
$\rho^{-s}$ has its only pole at $t=\rho$, with order
one more than its polynomial degree. Thus $P_\rho\ne0$
if and only if $\rho$ is a pole of $H_n$.
Let $L_0$ be the least common multiple of those root orders.
The second expression in (P-F56k) is periodic in $r$
with period $L_0$, so it exhibits $L_0$ as a quasipolynomial
period. Conversely $L_0\mid L$. Minimality gives $L=L_0$.
Now (P-F56a) and the primitive $M$th poles give (P-F56b).

In particular $L/M$ is odd. If $\rho$ is any Hilbert pole
other than a primitive $M$th root, its order has two-part
at most $M/2$ and divides $L/2$. Its contribution in the
difference of (P-F56k) at $r$ and $r+L/2$ therefore
vanishes. For $\rho=\zeta^a$ with $a$ odd,
$\rho^{-L/2}=-1$.
The simplicity and normalization calculation in P-F34 give
$P_{\zeta^a}(X)=A_n(\zeta^a)$, independent of $X$.
Subtracting the two Fourier formulas proves (P-F56d).
It also shows that every nonconstant coefficient of the
constituent polynomials has period dividing $L/2$.
The rational Hilbert series is defined over $\mathbb Q$,
so Galois conjugation sends $\mathcal C_n(\zeta)$ to
$\mathcal C_n(\zeta^a)$. This proves the trace expression
and rationality in (P-F56c). Its periodicity is immediate;
also $W_n(r+M/2)=-W_n(r)$.

After the coefficient-positivity shift $p=p(f_n)$, all
constituents evaluated at $p$ are nonnegative, since these
values are their shifted constant coefficients.
Take the residue $r\equiv p\pmod L$ in (P-F56d).
Its first term at $X=p$ is the actual value $f_n(p)$,
and its second term is nonnegative.
After a strict-saturation shift $s=s(f_n)$, every nonzero
constituent at $s+1$ is positive, and every identically
zero constituent has value zero. Taking
$r\equiv s+1\pmod L$ proves the second inequality of
(P-F56e). This uses exactly the source conventions, not
positivity only on each constituent's own progression.
These tests persist at every later integer argument.
For positivity, each polynomial $q_r(p+X)$ has nonnegative
coefficients and hence is nonnegative for all real $X\ge0$.
For saturation, the definition directly gives
$q_r(s+j)\ge0$ for every integer $j\ge1$.
Applying (P-F56d) at the residue of any such later argument
proves (P-F56e'). No derivative positivity is being inferred
from strict saturation.

Finally the multiplicity $f_n(x)$ is bounded by the
dimension of the ambient symmetric power of an
$m$-dimensional space. Since $m\le n^3$ and
$d_0x\le n^{K+1}$,
\[
 \binom{m+d_0x-1}{m-1}
      \le(m+d_0x)^{m-1}
      \le\bigl(2n^{\max\{3,K+1\}}\bigr)^{n^3},
\]
including $x=0$. This proves (P-F56f).
Its base-two logarithm is $O_K(n^3\log n)$.
The witness in (P-F56g) is at most $n^{K+1}$ for all
sufficiently large $n$. Apply (P-F56f) with $K+1$.
Its actual coefficient is smaller than $2W_n(x)$ for
sufficiently large $n\in\mathcal N$, so (P-F56e') gives
$p(f_n)>x\ge n^K$ and $s(f_n)\ge x\ge n^K$.
Applying this with a larger fixed exponent proves that
neither index is bounded by any polynomial in $n$.
The GCT parameter is $h=m^2+1+n=\Theta(n^6)$ by P-F31,
so a polynomial bound in that parameter would also fail.
This is a conditional deduction only; it does not prove
the wave estimate (P-F56g).\hfill$\square$
\end{proof}

\paragraph{The exact remaining task.}
The dyadic wave in (P-F56c) is an actual signed cyclotomic
trace of known simple-pole coefficients. Its control is
separate from the absolute amplitude estimate of P-F34.
This reduction cancels all other Hilbert-pole contributions
before any estimate is made. It does not assume that an
arbitrary affine slice inherits a radial period, that the
trace has a prescribed sign, or that local determinant
magnitudes cannot cancel. The sufficient wave condition
is not asserted to be equivalent to PH2 or SH, or logically
equivalent to the earlier average-polynomial sufficient gate.
The source index definitions are at
\url{https://arxiv.org/pdf/0704.0229},
Definitions~1.2.2 and 1.2.4, printed pp.~13--14.

\paragraph{A concrete dominant-pair sufficient condition.}
The amplitude needed for a bounded-window test can be located at a
specified pair of primitive roots. Put
\[
 \begin{gathered}
 r=3g\in\{3,9\},\qquad
 1\le a_*<M,\quad ra_*\equiv1\pmod M,\\
 Z_n=A_n(\zeta^{a_*}),\qquad
 E_n=\sum_{\substack{a\bmod M,\ a\ {\rm odd}\\
                         a\not\equiv a_*,-a_*\pmod M}}
                              |A_n(\zeta^a)|.
 \end{gathered}
 \tag{P-F56l}
\]
Here $Z_n$ is the actual normalized Hilbert residue, so its phase
is fixed by (P-F56c), not by a choice of tangent-versus-cotangent
convention. Unconditionally, for every $n\ge4$,
\[
 \begin{split}
 \log_2|Z_n|
 &\ge \binom n4-2\binom{n-1}3-m+1-(n-3)n(n-1)\\
 &\ge \frac{n^4}{24}-2n^3.
 \end{split}
 \tag{P-F56m}
\]
The following explicit remainder estimate is a sufficient
additional hypothesis; it is not proved here:
\[
                         E_n\le\frac14|Z_n|.
 \tag{P-F56n}
\]
For every $n\ge6$ for which (P-F56n) holds, and for every integer
$X$, there is an integer $x$ such that
\[
 X\le x\le X+r-1\le X+8,\qquad
 W_n(x)\ge\frac14|Z_n|
             \ge 2^{\,n^4/24-2n^3-2}.
 \tag{P-F56o}
\]
In particular, if (P-F56n) holds on an unbounded set
$\mathcal N$ of dimensions, then for every fixed real
$0<c<1/4$ and every sufficiently large $n\in\mathcal N$,
\[
                      p(f_n)>2^{cn},
             \qquad  s(f_n)\ge2^{cn}.
 \tag{P-F56p}
\]
Thus a specified total-remainder bound would give exponential
strict-index lower bounds. Neither (P-F56n) nor its index
consequence is asserted unconditionally.

\emph{Proof of the amplitude bound.}
Use representatives
$w_i=(-2)^i$ for $0\le i\le n-2$ and $w_{n-1}=0$,
and put $\eta=\zeta^{ga_*}$, so that $\eta^3=\zeta$.
The absolute determinant identity (P-F34n), which is unchanged
by inverting every eigenvalue, gives
\[
 |Z_n|=\frac1M
 \frac{\displaystyle\prod_{i\ne j}|1-\eta^{w_i-w_j}|}
      {\displaystyle\prod_{\substack{0\le i\le j\le k<n\\
                     w_i+w_j+w_k\not\equiv0\pmod M}}
                        |1-\eta^{w_i+w_j+w_k}|}.
 \tag{P-F56q}
\]
The $n$ variable weights are distinct modulo $M$.
Since $\eta$ is primitive, every factor in the numerator
is at least $4/M$: for $1\le u\le M/2$,
$2\sin(\pi u/M)\ge4u/M$.
There are $n(n-1)$ such factors.

Consider the $\binom{n-1}3$ distinct triples
$0\le i<j<k\le n-2$ that use no zero-weight variable.
Their integer weight
$w=w_i+w_j+w_k$ is divisible by three, because each
$(-2)^i$ is congruent to one modulo three.
It is nonzero modulo $M$, since the least two-adic
valuation occurs exactly once.
Also $|w|<2^{k+1}$. Consequently
\[
 |1-\eta^w|
   =|1-\zeta^{w/3}|
   \le\frac{2\pi|w|}{3M}
   \le\frac{8\,2^k}{M}.
 \tag{P-F56r}
\]
The last, deliberately loose constant follows from $\pi<4$.
The sum of $n-1-k$ over these triples is $\binom n4$:
it counts the choices of a fourth index larger than $k$.
Their reciprocal denominator product is therefore at least
$2^{\binom n4-3\binom{n-1}3}$.
There are exactly $m-n$ moving cubic monomials, since the
fixed ones are the $n$ chain monomials of P-F26.
Every remaining denominator factor is at most two.
Combining these facts with the numerator lower bound and
the normalization $1/M$ proves the first line of (P-F56m).
For the second line, use, for $n\ge4$,
\[
 \binom n4\ge n^4/24-n^3/4,\quad
 2\binom{n-1}3\le n^3/3,\quad
 m\le n^3/3,\quad
 (n-3)n(n-1)\le n^3.
 \]
Their total cubic loss is at most $23n^3/12<2n^3$.

\emph{Proof of the bounded-window and index implications.}
Write $ra_*=1+kM$. Then $1\le k<r$ and
$\gcd(k,r)=1$, since $kM\equiv-1\pmod r$.
Thus
$a_*/M=k/r+1/(rM)$.
For any fixed integer $X$, the $r$ complex numbers
\[
 Z_n\zeta^{-a_*X}\exp(-2\pi i k j/r),
                   \qquad 0\le j<r,
 \]
are equally spaced on the circle of radius $|Z_n|$.
At least one has real part at least
$\cos(\pi/r)|Z_n|\ge |Z_n|/2$.
Replacing its phase by the actual phase
$\zeta^{-a_*j}$ changes its real part by at most
$(2\pi/M)|Z_n|\le |Z_n|/4$ when $n\ge6$.
For this choice of $j$, therefore,
\[
 2\operatorname{Re}
        \bigl(Z_n\zeta^{-a_*(X+j)}\bigr)
                         \ge\frac12|Z_n|.
 \tag{P-F56s}
\]
The two terms with $a=a_*$ and $a=-a_*$
in (P-F56c) are precisely this conjugate pair:
complex conjugation preserves the rational Hilbert series.
All the remaining terms have total absolute value at most
$E_n$. Subtracting that error and using (P-F56n) proves
(P-F56o), independently of the initial phase of $Z_n$.

Finally take $X=\lceil2^{cn}\rceil$ with fixed $0<c<1/4$.
For the witness $x\le X+8$, the same ambient symmetric-power
bound used in (P-F56f), keeping $m$ explicit, gives
\[
 \log_2\max\{1,f_n(x)\}
 \le (m-1)\log_2(m+nx)
 =\frac c6 n^4+O_c(n^3\log n).
 \tag{P-F56t}
\]
Since $c/6<1/24$, this is eventually smaller than
$\log_2(2W_n(x))\ge n^4/24-2n^3-1$.
The one-witness implication (P-F56e') gives
$p(f_n)>x$ and $s(f_n)\ge x$, proving (P-F56p).
All estimates before (P-F56n) and the implications from it
are proved; the unsolved part is the single explicit
remainder estimate (P-F56n).\hfill$\square$

\paragraph{An exact positive multiset for the chain normal representation.}
There is also a numerator-free product specifying every residue
in the remaining estimate. Regard
$w_0,\ldots,w_{n-1}$ as residues modulo $M$, with the
representatives used in (P-F56q), and form the multiset
\[
 \begin{gathered}
 \mathcal E_n=
 \{\{0,1,k\}:2\le k<n\}\cup\{\{0,2,3\}\},\\
 \mathcal S_n=\{w_i:0\le i\le n-2\}\ \sqcup\
 \{w_i+w_j+w_k:0\le i<j<k<n,\
       \{i,j,k\}\notin\mathcal E_n\}.
 \end{gathered}
 \tag{P-F56u}
\]
Repeated weights are retained. Every weight is nonzero modulo
$M$, and $|\mathcal S_n|=D=\binom n3$.
With the actual Hilbert-residue convention, the exact identity is
\[
 \mathcal C_n(\zeta)=
                 \prod_{b\in\mathcal S_n}(1-\zeta^{gb})^{-1}.
 \tag{P-F56v}
\]
This is an equality of complex numbers, including their phases,
not only an absolute-value identity.

To prove it, put $\eta=\zeta^g$ and let
$\alpha_i=\eta^{w_i}$ be the natural eigenvalues of the diagonal
chain automorphism. Write $V$ for their representation.
Its action on polynomial coordinate functions is the usual
$GL_n$ action on $\operatorname{Sym}^3 V$.
On an invariant of coefficient degree $d_0s$ the character is
$(\det)^{\gamma s}$, where $\gamma=3/g$.
Since $3\sum_i w_i\equiv1\pmod M$, one has
$(\prod_i\alpha_i)^\gamma=\zeta$.
Thus its action on the canonical multiplicity algebra is precisely
the ring automorphism multiplying degree $s$ by $\zeta^s$.
This choice fixes the convention required in (P-F56v).

The smooth quotient tangent sequence in P-F24, dualized to
coordinate functions, shows that the moving conormal character is
\[
 \operatorname{Sym}^3 V-V\otimes V^*.
 \tag{P-F56w}
\]
Indeed the removed orbit cotangent character is the adjoint
character $V\otimes V^*-\mathbf1$, and the remaining fixed
grading direction contributes the other $\mathbf1$.
This is an identity in the representation ring of the finite
cyclic group, with its left side realized by the actual
conormal representation.

Let $h_j$ and $e_j$ be the complete and elementary symmetric
characters in the $\alpha_i$.
The identity
$\prod_i(1+t\alpha_i)=
 \prod_i(1-t^2\alpha_i^2)/\prod_i(1-t\alpha_i)$
gives $h_3=e_3+h_1\sum_i\alpha_i^2$.
The multisets $\{\alpha_i^2\}$ and $\{\alpha_i^{-1}\}$
differ by replacing $\alpha_0^{-1}$ with one: this follows
from $\alpha_i^2=\alpha_{i+1}^{-1}$ for $i<n-1$
and $\alpha_{n-1}=1$.
Consequently (P-F56w) is
\[
                       e_3+h_1-\alpha_0^{-1}h_1.
 \tag{P-F56x}
\]
The trivial weight in $h_1$ cancels the term with $i=0$
in $\alpha_0^{-1}h_1$.
For every $2\le k<n$, the remaining subtracted weight
$w_k-w_0=w_k-1$ equals the weight of the triple
$\{0,1,k\}$, because $w_0+w_1=-1$.
The last subtracted weight is $w_1-w_0=-3$.
It equals the weight of $\{0,2,3\}$:
for $n\ge5$ the integer sum is $1+4-8=-3$;
for $n=4$ it is $1+4+0=5\equiv-3\pmod8$.
These removed triples are distinct.
The remaining character is exactly the positive multiset
(P-F56u). The count is
$(n-1)+\binom n3-(n-2)-1=D$.
Nonvanishing of its singleton weights is immediate, and
a distinct triple sum is nonzero by its unique smallest
two-adic valuation among the nonzero summands.
The residue formula (P-F31b) now gives (P-F56v).

For reference, put
$a_0=g(n+1)(n+2)/6\in\mathbb Z$.
The determinant of $\operatorname{Sym}^3 V$ is
$(\det V)^{\binom{n+2}2}$, whereas
$\det(V\otimes V^*)=1$.
It follows from (P-F56w) and $3\sum_i w_i\equiv1\pmod M$
that the sum of the exponents $gb$ in (P-F56v) is congruent
to $a_0$ modulo $M$.
Inverting each factor therefore gives the exact relations
\[
 \mathcal C_n(\zeta^{-1})
       =(-1)^D\zeta^{a_0}\mathcal C_n(\zeta),
 \qquad
 W_n(s)=(-1)^D W_n(-s-a_0).
 \tag{P-F56y}
\]
They fix a reciprocity symmetry, but do not prescribe the sign
of the wave at a specified nonnegative integer.

Finally, if $h=M/2$ and
$\mathcal C_n(\zeta)=\sum_{j=0}^{h-1}c_{n,j}\zeta^j$
is its unique representative in
$\mathbb Q[t]/(t^h+1)$, then for every integer $s$,
\[
             2W_n(s)=
       (-1)^{\lfloor s/h\rfloor}c_{n,\,s\bmod h}.
 \tag{P-F56z}
\]
Indeed the primitive dyadic trace of $\zeta^b$ is zero unless
$h\mid b$, and is $h(-1)^{b/h}$ otherwise. This follows
directly by summing the geometric series over the odd
exponents modulo $M$.
Only $j=s\bmod h$ survives in (P-F56c), giving (P-F56z).
In particular, at $0\le x<h$ the exact one-witness
obstruction is $f_n(x)<c_{n,x}$.
This is a coefficient reformulation of the unresolved signed
wave question, with no positivity of those coefficients assumed.

\paragraph{A contraction gate involving only quadratic-many factors.}
The unresolved mass comparison can be reduced further to a
dimension-increment inequality. This reduction proves the error
from the two selected sibling branches; it does not assume or
prove the remaining uniform contraction inequality.

For $n\ge4$ write $M_n=2^{n-1}$, let
$b_n=3^{-1}\pmod {M_n}$, and use the positive multiset
$\mathcal S_n$ of (P-F56u). For every odd residue $a\bmod M_n$
define
\[
 P_n(a)=
 \prod_{b\in\mathcal S_n}
       |1-\exp(2\pi iab/M_n)|^{-1},\qquad
 Q_n=\sum_{\substack{a\bmod M_n,\ a\ {\rm odd}\\
                               a\ne b_n,-b_n}}
                              \frac{P_n(a)}{P_n(b_n)}.
 \tag{P-F56aa}
\]
Every factor is nonzero. Multiplication of the primitive-root
index by the odd number $g=\gcd(n,3)$ permutes the odd residues.
Thus (P-F56v) implies exactly $Q_n=E_n/|Z_n|$; the common
normalization $1/M_n$ cancels.

For $n\ge5$, form the multiset of residues modulo $M_n$
\[
 \begin{split}
 \mathcal T_n={}&\{1\}\\
 &\sqcup\{1+w_j+w_k:2\le j<k<n,\ (j,k)\ne(2,3)\}\\
 &\sqcup\{2+w_k:3\le k<n\}\sqcup\{6\},
 \qquad |\mathcal T_n|=\binom{n-1}2,
 \end{split}
 \tag{P-F56ab}
\]
where $w_i=(-2)^i$ for $i<n-1$ and $w_{n-1}=0$,
as before. Define the positive increment product
$J_n(a)=\prod_{b\in\mathcal T_n}
 |1-\exp(2\pi iab/M_n)|^{-1}$.
The exact recursion is
\[
             P_n(a)=J_n(a)P_{n-1}(-a\bmod M_{n-1}).
 \tag{P-F56ac}
\]

Here is a sufficient uniform contraction hypothesis, involving
only the $\binom{n-1}2=O(n^2)$ factors of each increment:
there exist a fixed $0<\rho<1$ and an integer $n_0\ge5$ such that,
for every $n\ge n_0$ and every odd $a\bmod M_n$ with
$a\bmod M_{n-1}\notin\{b_{n-1},-b_{n-1}\}$,
\[
 J_n(a)+J_n(a+M_{n-1})\le \rho J_n(b_n).
 \tag{P-F56ad}
\]
This inequality remains unproved. If it holds, then
$Q_n\longrightarrow0$, hence (P-F56n) holds for every
sufficiently large $n$ and gives the conditional exponential
index conclusion (P-F56p).
In this implication there is no unproved finite base
certificate: any finite starting value of $Q_n$ suffices.

The selected sibling error is bounded unconditionally:
\[
 \delta_n:=
 \frac{P_n(b_n+M_{n-1})}{P_n(b_n)}
 \le
 2^{-n^3/6+4n^2-107n/6+19},
 \qquad
 \delta_n\le2^{-n^3/12}\quad(n\ge48).
 \tag{P-F56ae}
\]
The exact implication of (P-F56ad) is
\[
                         Q_n\le\rho Q_{n-1}+2\delta_n.
 \tag{P-F56af}
\]

\emph{Proof.}
Remove the first variable from the chain weights.
The remaining weights $w_1,\ldots,w_{n-1}$ are $-2$
times the weights of the $(n-1)$-variable chain, in the
same order. In the multiset (P-F56u), all triples with
smallest index at least one occur. The shifted smaller
multiset removes from these exactly the triples
$\{1,2,k\}$ for $3\le k<n$ and $\{1,3,4\}$.
The additional triples with smallest index zero are
$\{0,j,k\}$ for $2\le j<k<n$, except $\{0,2,3\}$.
The additional singleton is $w_0=1$.
These extra weights are (P-F56ab): the first removed
family has weight $w_1+w_2+w_k=2+w_k$, and the last
triple has weight six modulo $M_n$.
For $n\ge6$ its integer weight is $-2-8+16=6$;
for $n=5$ it is $-2-8+0=-10\equiv6\pmod {16}$.
Thus the multiset decomposition is exact, including
multiplicities and the smallest allowed dimension.
Its cardinality is
$1+(\binom{n-2}2-1)+(n-3)+1=\binom{n-1}2$.
All its factors are nonzero because it is a submultiset
of $\mathcal S_n$. Substitution of the root
$\exp(-4\pi ia/M_n)$ in the smaller product proves
(P-F56ac).

Put $M=M_n$ and $\eta=\exp(2\pi ib_n/M)$.
The two roots indexed by $b_n$ and $b_n+M/2$
have the same smaller-dimensional factor in (P-F56ac).
In the increment, every even weight gives the same
factor at $\eta$ and $-\eta$.
The odd weights of $\mathcal T_n$ are exactly
\[
 \mathcal O_n=
 \{1\}\sqcup
 \{1+w_j+w_k:2\le j<k<n,\ (j,k)\ne(2,3)\}.
 \]
It follows that
\[
                 \delta_n
       =\prod_{b\in\mathcal O_n}
                     \left|\tan\frac{\pi b_n b}{M}\right|.
 \tag{P-F56ag}
\]
No tangent is singular, since $b_n$ and every $b$ here
are odd and $M\ge16$.

Among these odd weights, the bulk pairs
$2\le j<k\le n-2$, excluding $(2,3)$,
give
$B=\binom{n-3}2-1$ weights divisible by three.
For such a weight $b=1+w_j+w_k$, one has
$|b|<2^{k+1}$ and $|b|\le M$.
Since $3b_n\equiv1\pmod M$,
$\tan(\pi b_n b/M)=\tan(\pi b/(3M))$.
On $[-\pi/3,\pi/3]$, the bounds
$|\sin u|\le|u|$ and $\cos u\ge1/2$ give
\[
 \left|\tan\frac{\pi b_n b}{M}\right|
       \le\frac{2\pi |b|}{3M}
       \le\frac{8\,2^k}{M}.
 \tag{P-F56ah}
\]
The sum of $n-1-k$ over these bulk pairs is
$\binom{n-2}3-(n-4)$: before the one excluded pair,
it counts triples $2\le j<k<\ell\le n-1$.
There are $n-2$ exceptional odd factors, namely the
singleton one and the pairs using the zero-weight
variable. Each is at most $M/2$.
Indeed an odd residue modulo $M$ is at distance at least
one from $M/2$, so its cosine in (P-F56ag) has absolute
value at least $\sin(\pi/M)\ge2/M$, while its sine has
absolute value at most one.
Combining these estimates gives
\[
 \begin{split}
 \log_2\delta_n
 &\le 3\left(\binom{n-3}2-1\right)
        -\binom{n-2}3+(n-4)+(n-2)^2\\
 &=-n^3/6+4n^2-107n/6+19 .
 \end{split}
 \tag{P-F56ai}
\]
For $n\ge48$, $4n^2\le n^3/12$ and
$-107n/6+19<0$, proving the last assertion of
(P-F56ae).

Finally group the odd residues modulo $M_n$ into their
two lifts above an odd residue modulo $M_{n-1}$.
Reduction of $b_n$ is $b_{n-1}$.
The two selected parent classes therefore have the two
selected children $\pm b_n$ and the two unselected
children $\pm b_n+M_{n-1}$.
Complex conjugation makes the latter two contributions
equal, with total normalized mass $2\delta_n$.
For every other parent, (P-F56ac) and (P-F56ad)
bound its two children by $\rho$ times its normalized
mass in $Q_{n-1}$.
Here
$P_n(b_n)=J_n(b_n)P_{n-1}(b_{n-1})$,
using conjugation for the minus sign in (P-F56ac).
This proves (P-F56af), with no factor of two lost
in the nonselected branches.

Since $\delta_n\to0$ and $\rho<1$, the recursion forces
$Q_n\to0$. Explicitly, for any $\epsilon>0$, choose
$N$ so that $2\delta_n\le\epsilon(1-\rho)/2$
for all $n>N$. Iterating gives
\[
 Q_n\le \rho^{n-N}Q_N+\epsilon/2
                  \quad(n\ge N),
 \tag{P-F56aj}
\]
and the first term tends to zero.
The starting value $Q_N$ is finite by its definition.
This completes the conditional contraction reduction
and the unconditional selected-sibling bound.

\paragraph{Uniform comparison on polynomial ranges of dyadic multipliers.}
The positive multiset (P-F56u) gives a rigorous part of the
remaining estimate (P-F56n). This does not bound the entire
remainder. Let $n\ge5$, $l=n-1$, $M=2^l$, and let $J$ be an
odd integer with $1\le J<M$. Put
\[
 a_J\equiv J(3g)^{-1}\pmod M,\qquad
 Z_{n,J}=A_n(\zeta^{a_J}),\qquad
 R_J=\lceil\log_2(16J)\rceil.
 \tag{P-F56ak}
\]
For negative subscripts use $a_{-J}\equiv-a_J\pmod M$;
the rational Hilbert series gives
$Z_{n,-J}=\overline{Z_{n,J}}$.
Then the following explicit uniform error bound holds:
\[
 \left|\log_2\frac{|Z_{n,J}|}{|Z_{n,1}|}
       +\binom{n-1}{3}\log_2J
       -\mathbf1_{3\mid J}\binom n3\right|
       \le 100n^2R_J^2.
 \tag{P-F56al}
\]
In particular, for every fixed real $K\ge0$, uniformly for
odd $3\le J\le n^K$,
\[
 \log_2\frac{|Z_{n,J}|}{|Z_{n,1}|}
 \le -\frac{\log_2 3-1}{6}n^3
                         +O_K(n^2\log^2 n).
 \tag{P-F56am}
\]
Consequently the sum of the magnitudes of all terms indexed
by $a\equiv\pm J(3g)^{-1}$ with $3\le J\le n^K$ is
$o(|Z_{n,1}|)$. Thus every fixed polynomial range of
multipliers is negligible in (P-F56n). This leaves the
remaining, unrestricted multipliers uncontrolled.

\emph{Sparse dyadic separation.}
We first give the elementary estimate that prevents a high
dyadic block from being cancelled by a very small tail.
Let $1\le r\le3$, $1\le q_1<\cdots<q_r$ be integers,
$\epsilon_i\in\{1,-1\}$, and $\kappa\in\mathbb Z$.
For odd $J\ge1$ and $R=\lceil\log_2(16J)\rceil$, set
\[
 x=\frac\kappa3+
              \frac J3\sum_{i=1}^r\epsilon_i2^{-q_i}.
\]
Then
\[
 \|x\|_{\mathbb R/\mathbb Z}
       \ge\frac1{6\,2^{q_1+(r-1)R}}.
 \tag{P-F56an}
\]
To prove this, retain the initial block ending just before
the first gap $q_{i+1}-q_i>R$, or retain all terms if no
such gap occurs. If the retained block ends at $q_t$, its
numerator over the denominator $3\cdot2^{q_t}$ is odd:
the final term is odd, all earlier terms are even, and
$\kappa2^{q_t}$ is even. It is therefore not an integer
and its distance from an integer is at least
$1/(3\cdot2^{q_t})$.
If a tail remains, its absolute value is less than
\[
 \frac{2J}{3}\,2^{-q_{t+1}}
 \le\frac{J2^{-R}}{3\,2^{q_t}}
 \le\frac1{48\,2^{q_t}}.
\]
The reverse triangle inequality for distance to the
integers now gives at least $1/(6\cdot2^{q_t})$.
Since every retained gap is at most $R$, one has
$q_t\le q_1+(r-1)R$, proving (P-F56an).
In particular, if $q_1\le R$ and $r\le3$, then
\[
 2^{-3R-1}\le |1-e^{2\pi ix}|\le2.
 \tag{P-F56ao}
\]
Here we used $|1-e^{2\pi ix}|\ge4\|x\|$.

\emph{Partition of the exact product.}
Write $w_i=(-2)^i$ for $0\le i<l$, with the remaining
weight zero. Let $b_J$ be the residue $J3^{-1}$ modulo
$M$. Then $\eta_J=\zeta^{b_J}$ satisfies
$\eta_J^3=\zeta^J$, and
\[
 \eta_J=\omega_J\exp\left(\frac{2\pi iJ}{3M}\right),
 \qquad \omega_J^3=1,
 \qquad \omega_J=1\ \Longleftrightarrow\ 3\mid J.
 \tag{P-F56ap}
\]
Indeed $3b_J=J+k_JM$ with $k_J\in\{0,1,2\}$, and
$\omega_J=e^{2\pi i k_J/3}$; because $J<M$, the only
case $k_J=0$ is $3\mid J$.
Every $w_i$ is congruent to one modulo three. The exact
denominator in (P-F56v) therefore splits into:
\begin{enumerate}
\item the $l$ singleton factors $1-\eta_J^{w_i}$;
\item the distinct active triple factors
      $1-\eta_J^{w_i+w_j+w_k}$, $i<j<k<l$,
      omitting $\{0,1,k\}$ and $\{0,2,3\}$;
\item the factors $1-\eta_J^{w_i+w_k}$,
      $0\le i<k<l$, omitting the pair $\{0,1\}$.
\end{enumerate}
The last family consists precisely of triples using the
zero-weight variable. The common factor $1/M$ in the
normalized residues cancels from their ratio.

For any factor let $k$ be the largest active index and put
$q=l-k$. Call it low if $q>R$ and high if $q\le R$,
where $R=R_J$. A sum of at most three distinct active
weights divided by $M$ has the form
$\sum\epsilon_i2^{-q_i}$ with smallest denominator
exponent $q_1=q$. Its absolute value lies between
$2^{-q}/4$ and $2\cdot2^{-q}$ when it has three terms;
for two terms the sharper bounds are
$2^{-q}/2$ and $3\cdot2^{-q}/2$.
All phases are consequently covered by (P-F56an).

\emph{Active triples.}
For an active triple, its integer weight is divisible
by three, so its phase at $J$ is $J$ times its phase
at $1$, with no cube-root shift. On the low factors
the absolute phase is at most
$(2J/3)2^{-q}<1/24$.
The elementary inequalities
$1-y^2/6\le\sin y/y\le1$ show, for $|x|\le1/24$,
\[
 \left|\log_2\frac{\sin(\pi x)}{\pi x}\right|
                                      \le8x^2.
\]
Thus the logarithm of the ratio of the $J$ factor to
the $1$ factor differs from $\log_2 J$ by at most
$16J^2 4^{-q}$. For each $q$ there are at most
$n^2/2$ triples. Summing the geometric series over
$q>R$ gives total error less than $n^2$, since
$2^R\ge16J$.

There are at most $Rn^2/2$ high triples.
Their $J$ factors satisfy (P-F56ao). Their reference
factors have absolute value at least $2^{-q}/3$:
the reference phase has absolute value between
$2^{-q}/12$ and $2\cdot2^{-q}/3$, is less than $1/2$,
and $|1-e^{2\pi ix}|\ge4|x|$ on this interval.
The absolute logarithm of each factor ratio is therefore
at most $5R$. Replacing the number of low triples by
$\binom l3$ costs at most another $Rn^2\log_2 J/2$.
The omitted active triples number $l-1$; each ratio
has absolute logarithm at most $5R$ (or at most $2R$
in the low case). Combining these bounds, their net
contribution to the logarithm of the residue ratio is
\[
               -\binom l3\log_2 J+E_3,
                         \qquad |E_3|\le10n^2R^2.
 \tag{P-F56aq}
\]
The loose constant also covers $R\ge l$, when there
are no low factors.

\emph{Pairs using the zero-weight variable.}
At the reference embedding the phase of a pair is
$1/3$ or $2/3$, plus a perturbation of absolute value
at most $1/4$. Hence every reference factor has
absolute value between $1/3$ and $2$.
If $3\nmid J$, the low target factors also have a
nontrivial cube-root shift, perturbed by less than
$1/24$; they and the reference factors have logarithms
bounded by absolute constants.
There are at most $nR$ high pairs. Applying the
two-term version of (P-F56an) bounds the absolute
logarithm of each high ratio by $3R$.
Their total contribution is therefore bounded by
$4n^2R^2$ in absolute value.

If $3\mid J$, the low target phase has no cube-root
shift. Its factor lies between
$(2/3)J2^{-q}$ and $4J2^{-q}$.
The logarithm of the reciprocal-factor ratio is
therefore $q-\log_2 J$ with error at most four.
For the full pair family,
\[
 \sum_{0\le i<k<l}(l-k)
      =\sum_{k=1}^{l-1}k(l-k)
      =\binom{l+1}{3}=\binom n3.
\]
Discarding the $-\log_2 J$ terms costs at most
$n^2\log_2 J/2$. The high pairs number at most $nR$,
have $q\le R$, and their logarithmic ratios are at
most $3R$ in absolute value. Omitting the pair
$\{0,1\}$ costs at most $n+\log_2 J+4$ in the low
case, or $3R$ in the high case. All these bounds give
the contribution
\[
           \mathbf1_{3\mid J}\binom n3+E_2,
                         \qquad |E_2|\le8n^2R^2.
 \tag{P-F56ar}
\]

\emph{Singletons and conclusion.}
Reference singleton factors have absolute value at
least $2/3$. If $3\nmid J$, the low target factors
stay a fixed distance from one; the at most $R$
high factors satisfy the one-term separation estimate.
If $3\mid J$, each low logarithmic ratio has absolute
value at most $q+\log_2 J+3$, and the high ratios have
absolute value at most $R+2$.
Summing over the $l$ singletons in either case gives
absolute contribution at most $5n^2R^2$.
Together with (P-F56aq)--(P-F56ar), this proves
(P-F56al), with room in the stated constant $100$.

For $J\le n^K$, $R=O_K(1+\log n)$.
Since $J\ge3$, the main terms in (P-F56al) are at most
$-(\log_2 3-1)n^3/6+O_K(n^2\log n)$.
This proves (P-F56am). Complex conjugation gives equal
magnitudes for the multipliers $J$ and $-J$.
For all sufficiently large $n$ the range $J\le n^K$
is smaller than $M/2$, so these indices are distinct.
There are at most $2n^K$ terms, whence
\[
 \limsup_{n\to\infty}\frac1{n^3}\log_2
 \left(
 \frac{\displaystyle\sum_{\substack{3\le J\le n^K\\J\text{ odd}}}
       (|Z_{n,J}|+|Z_{n,-J}|)}{|Z_{n,1}|}
 \right)
 \le-\frac{\log_2 3-1}{6}<0.
 \tag{P-F56as}
\]
An empty sum is interpreted as zero, with logarithm
$-\infty$. This proves the claimed negligible part
of the actual remainder. No comparison of the other
multipliers, and hence no unconditional strict-index
lower bound, follows from this argument.
\hfill$\square$

\paragraph{An exponentially long range with a quantitative bound.}
The explicit error in (P-F56al) gives more than polynomial
ranges. For every odd integer $3\le J\le2^{n/100000}$,
\[
       \log_2\frac{|Z_{n,J}|}{|Z_{n,1}|}
                      \le-\frac{n^3\log_2 J}{36}.
 \tag{P-F56at}
\]
The range may be empty in small dimensions. In its stated
range it satisfies $J<M$ for every $n\ge5$.
For a proof put $t=\log_2J>3/2$. The elementary bounds
$\binom{n-1}{3}\ge n^3/6-n^2$ and
$\binom n3\le n^3/6$, followed by (P-F56al), give
\[
 \log_2\frac{|Z_{n,J}|}{|Z_{n,1}|}
 \le-\frac{(t-1)n^3}{6}
                      +n^2\bigl(t+100(t+5)^2\bigr).
\]
Here $R_J\le t+5$. Since $t-1\ge t/3$ and
$t+5\le5t$, the last bound is at most
$-tn^3/18+2501n^2t^2$.
Now $t\le n/100000$ and $2501/100000<1/36$ prove
(P-F56at). Summing and using conjugation yields the
fully explicit bound
\[
 \frac{\displaystyle\sum_{\substack{3\le J\le2^{n/100000}\\
                                  J\text{ odd}}}
       (|Z_{n,J}|+|Z_{n,-J}|)}{|Z_{n,1}|}
 \le 2^{\,1+n/100000-(\log_2 3)n^3/36}.
 \tag{P-F56au}
\]

There is also a sharper asymptotic range. For every fixed
$0<\delta<1/600$, the conclusion (P-F56as), with the
same negative upper bound, holds when $n^K$ is replaced
by $2^{\delta n}$. We spell out the uniformity argument.
Consider any sequence of admissible odd $J=J_n$ and put
$t_n=\log_2J_n$. If $(t_n)$ is bounded, pass to a
subsequence with the integer $J_n$ constant. Dividing
(P-F56al) by $n^3$ gives upper limit
$(-\log_2J+\mathbf1_{3\mid J})/6$, which is at most
$-(\log_2 3-1)/6$ for $J\ge3$.
If $(t_n)$ is unbounded, pass to a subsequence on which
$t_n\to\infty$. The displayed estimate above, divided
by $n^3$, is
\[
 -\frac{t_n-1}{6}
        +\frac{t_n+100(t_n+5)^2}{n}
 \le-\bigl(1/6-100\delta-o(1)\bigr)t_n+O(1),
\]
because $t_n\le\delta n$. This tends to $-\infty$.
Every hypothetical failure of the uniform upper limit
would have one of these two subsequences, so neither
is possible. Finally there are at most $2^{1+\delta n}$
conjugate terms, whose logarithm divided by $n^3$ tends
to zero. This proves the asserted uniform exponential-range
version of (P-F56as). The multipliers outside that range
remain an open part of (P-F56n).
\hfill$\square$
